\title{Large Language Models Can Automatically Engineer Features for Few-Shot Tabular Learning}

\begin{abstract}
Large Language Models (LLMs) possess significant potential to address tabular learning, a crucial component in numerous real-world scenarios, due to their extraordinary aptitude for solving difficult and novel reasoning challenges. In this study, we introduce \textsf{FeatLLM}, an innovative in-context learning strategy that utilizes LLMs as feature creators, constructing an input dataset optimized for predictive modeling with tabular data. The features generated through this approach are then leveraged by straightforward downstream machine learning algorithms, such as linear regression, enabling highly effective few-shot learning performance. Notably, \textsf{FeatLLM} relies solely on this basic predictor in conjunction with the newly constructed features for inference, eliminating the necessity to query the LLM for each individual input during test time as is typical in many LLM-based solutions. Furthermore, \textsf{FeatLLM} operates with just API access to LLMs and successfully bypasses constraints on prompt sizes. Comprehensive experiments over diverse tabular datasets demonstrate that \textsf{FeatLLM} yields superior rule generation, outperforming counterparts including TabLLM and STUNT by a margin of approximately 10\% on average.
\end{abstract}

\section{Introduction}
Large Language Models (LLMs) have recently made striking advances, demonstrating proficiency in generating accurate outputs for novel tasks without the need for extensive task-specific adjustment~\cite{creswell2022selection,imani2023mathprompter,taylor2022galactica}. Due to training on massive text datasets, LLMs inherit extensive world knowledge, often replacing traditional reliance on domain expertise~\cite{Genomes2021,singhal2023large,wilcox2003role}. This intrinsic knowledge base supports automating a spectrum of tasks, including but not limited to dialogue analysis~\cite{finch2023leveraging} and program repair~\cite{fan2023automated}. Particularly, by integrating descriptions of tasks and small exemplar sets within prompts, LLMs are capable of contextual comprehension—identifying salient examples and features~\cite{brown2020language,zhao2023survey}—thereby illustrating their potential as skilled feature engineers.

Applying the reasoning abilities and rich prior information of LLMs has become increasingly popular in the realm of tabular learning. As an example, knowledge such as higher disease risks among older adults can enhance disease prediction efforts. Recent methods translate tabular tasks and feature explanations into natural language formats, sequence them, and introduce these structured descriptions as inputs for LLMs~\cite{dinh2022lift,wang2023anypredict}. Approaches commonly employ in-context learning~\cite{nam2023semi} or parameter-efficient adaptation techniques~\cite{dinh2022lift,hegselmann2023tabllm} following this data preparation. Harnessing LLM-derived priors has been demonstrated to surpass traditional tabular learning models in settings with limited training samples~\cite{hegselmann2023tabllm}.

Existing tabular learning strategies based on LLMs exhibit several shortcomings. For methods that generate predictions directly through LLMs, an inference must be performed for each input sample, resulting in high computational overhead. Furthermore, achieving strong performance typically necessitates fine-tuning the LLM, which becomes impractical for models lacking open-source training code or available tuning interfaces—a growing issue as recent state-of-the-art LLMs often restrict access to inference APIs only~\cite{anil2023palm,openai2023gpt4}. Additionally, these techniques generally do not scale well with lengthy prompts~\cite{liu2023lost}; as the count of features in tabular datasets increases, and the textual representation grows accordingly, many such approaches become intractable or their effectiveness substantially degrades. Collectively, these issues limit the deployment of LLM-powered tabular learning solutions in practical environments.

We take a different direction by assigning LLMs to feature engineering rather than using them solely for end-to-end inference. Our method leverages LLMs' prior knowledge in a manner that is both more practical and efficient. Specifically, we present \textsf{FeatLLM}, an innovative in-context learning paradigm designed to decipher the implicit ``criteria" guiding LLM predictions. For example, in medical classification tasks, the LLM can be prompted to generate interpretable decision rules—such as the set of feature conditions indicative of a particular disease. These extracted criteria serve as newly engineered features, which can substitute the original feature set for downstream modeling. This refocuses the LLM’s contribution toward feature creation and allows for the use of simple models in subsequent stages, thereby reducing inference time and complexity. Utilizing in-context learning capabilities, our approach enables feature extraction by merely appending illustrative examples to the prompt, circumventing the need for any additional model training. To further address the issue of oversized prompts, we integrate ensemble-based strategies like feature bagging~\cite{breiman2001random}, which help manage prompt length efficiently.

\textsf{FeatLLM} decomposes the process of crafting new features via prompting into two distinct stages: first, it interprets the underlying problem to identify how features relate to target outputs; second, utilizing insights from this analysis alongside a handful of illustrative samples, it formulates specific rules for class differentiation. This sequential reasoning strategy guides the LLMs to disregard irrelevant features when creating rules. The resultant rules are then translated into executable code, which processes the data by converting compliance with each rule into a binary indicator. Subsequently, a simple model is employed to predict class membership probabilities from these generated features. To bolster reliability, an ensemble method is utilized, whereby feature generation and selection are performed multiple times through feature bagging. Adjusting the number of selected features and samples in each bag effectively manages prompt length, alleviating issues of prompt overflow. Since \textsf{FeatLLM} bypasses the need for explicit training, inference is cost-effective and the LLM is readily adaptable to a range of practical applications.

We assess \textsf{FeatLLM} across 13 tabular datasets under limited data scenarios and find it exhibits notable robustness and effectiveness. Results demonstrate that LLMs can draw on both domain knowledge and dataset-specific information to generate highly relevant features. Our method consistently surpasses state-of-the-art few-shot learning competitors in multiple contexts. The source code is made available through an anonymized GitHub repository at~\texttt{https://github.com/Sungwon-Han/FeatLLM}.

\section{Related Work}

\subsection{Few-Shot Learning with Tabular Data} Few-shot learning methods have facilitated the development of highly transferable representations from datasets where only limited labeled examples are available~\cite{chen2018closer}. While early research was centered on image-based tasks, few-shot approaches have increasingly demonstrated versatility across different modalities, such as text, audio, and tabular formats~\cite{majumder2022few,nam2023stunt,schick2021exploiting}. The challenge inherent to learning from sparse labeled data is the heightened susceptibility to capturing incidental correlations~\cite{bartlett2021deep}. To mitigate this, contemporary techniques incorporate supplementary information or domain-relevant priors to imbue models with more suitable inductive biases. This can involve strategies like exploiting large pools of unlabeled samples, augmented thoughtfully for semi-supervised paradigms~\cite{ucar2021subtab,yoon2020vime}, or employing generic unsupervised pre-training for improved model initialization~\cite{somepalli2021saint}. Unsupervised methods often mask or perturb segments of the dataset and reconstruct them~\cite{arik2021tabnet,majmundar2022met}, or apply augmentation schemes conducive to contrastive learning objectives~\cite{bahri2022scarf,chen2020simple}. Nevertheless, due to the heterogeneous characteristics of tabular data, universally effective augmentation methods remain elusive, and existing augmentation-based techniques have shown limited efficacy in few-shot regimes~\cite{nam2023stunt}.

A more recent trend has been to train models on an assortment of tabular datasets not confined to the target task, followed by transferring learned knowledge to the target application~\cite{zhang2023generative,levin2022transfer,zhu2023xtab}. For example, TransTab~\cite{wang2022transtab} leverages transformer architectures to acquire semantic insights between table columns by utilizing column names in addition to cell values. This cross-table knowledge transfer has enabled models to undertake zero-shot or few-shot prediction tasks on previously unseen tabular problems. Similarly, TabPFN~\cite{hollmann2022tabpfn} utilizes mixtures of statistical distributions to synthetically generate pre-training data for Bayesian neural network models. The pre-trained model infers from both training and test examples of the target task with no further training required. The prior accumulated from varied datasets has yielded performance superior to classic tree-based techniques and has proven advantageous for few-shot tabular prediction.

\subsection{Language-Interfaced Tabular Learning} Incorporating large language models (LLMs)~\cite{anil2023palm,openai2023gpt4} stands out as another avenue for harnessing prior knowledge in tabular tasks. With their training on vast, multifaceted text sources, LLMs exhibit remarkable generalization to novel tasks and are applicable across numerous domains~\cite{brown2020language}. Recent investigations have sought to leverage the broad prior knowledge embedded within LLMs for tabular reasoning. Approaches commonly involve transforming tabular information into textual form and supplying these as prompts to the LLMs~\cite{dinh2022lift,hegselmann2023tabllm,wang2023anypredict}. Incorporation of tabular entries, feature annotations, and task explanations in the input enables models to deduce task-feature associations for inference. Development has progressed towards specialized adaptation using methods like LoRA~\cite{hu2021lora} or IA3~\cite{liu2022few}, which allow parameter-efficient fine-tuning, or through in-context techniques where few-shot exemplar demonstrations are embedded within prompts without further model training~\cite{dinh2022lift,hegselmann2023tabllm,nam2023semi,slack2023tablet}.

Although the methods previously described have demonstrated efficacy in advancing few-shot tabular learning, their reliance on continual inference using an LLM introduces obstacles for deployment in industry settings. This research seeks to shift away from the traditional paradigm of using LLMs as end-to-end black-box predictors for each sample. Rather, the objective is to have the LLM directly identify the conditions or rules pertinent to each class. \textsf{FeatLLM} translates these generated rules into programmatic code that produces new features, enabling inference without the LLM and capitalizing on in-context learning to derive rules from prior knowledge and the few-shot data.

\section{Methods}
\textsf{FeatLLM} operates by extracting ``rules"—that is, the specific conditions characterizing each answer class—from a limited number of input samples, as illustrated in Figure~\ref{fig:main_model}. These rules are obtained using an LLM and are subsequently converted to binary features for each sample, denoting whether the sample meets each particular rule. The resulting collection of binary features is fed into a dot product with non-negative trainable parameters, yielding an estimate of class probability. Through model training, the significance of each feature in determining class membership is learned (Section~\ref{Sec:training}). This procedure is performed repeatedly in a bagging manner, with final predictions derived by ensemble aggregation. The following subsections outline the formal problem setup and elaborate on each operational phase.

\vspace*{-0.12in}
\paragraph{Problem formulation.} Consider a tabular dataset comprising $N$ labeled examples, denoted as $\mathcal{D} = \{(\mathbf{x}^i, \mathbf{y}^i)\}_{i=1}^{N}$. Each sample $\mathbf{x}^i$ contains $d$ input features, forming a $d$-dimensional vector. The dataset $\mathcal{D}$ is accompanied by feature names expressed in natural language, $F = \{f_j\}_{j=1}^{d}$, and each feature's definition is available separately. In the context of classification, each label $\mathbf{y}^i$ represents one of a finite set of class outcomes. For the $k$-shot learning scenario, only a subset of $k$ ($<N$) labeled instances is randomly selected for model training.

\subsection{Prompt Design for Extracting Rules}
\label{Sec:prompt_design}
To facilitate the extraction of rules from an LLM by guiding its reasoning with higher fidelity, we structure the prompting methodology to reflect how an expert would tackle a tabular learning problem. Our prompt consists of three primary elements (illustrated in Fig.~\ref{fig:prompt_for_rule} and detailed in Appendix~\ref{sec:example_rule}):

\vspace*{-0.12in}
\paragraph{Basic information description.} This section provides core details necessary for problem-solving (refer to orange highlights in Fig.~\ref{fig:prompt_for_rule}). It encompasses the task statement—including label information—as well as descriptions and definitions of the features, plus representative examples. Tasks are posed as questions, such as: ``Does this patient's myocardial infarction complication data suggest chronic heart failure? Yes or no?". Further task description instances can be found in Appendix~\ref{sec:task_desc}. Feature descriptions clarify each value type and may supply additional definitions; for categorical attributes, possible value categories are optionally listed. For exemplification, several training instances are serialized into textual format along with their true labels. The serialization process adheres to conventions established in prior works~\cite{dinh2022lift,hegselmann2023tabllm}:

\begin{align}
&\text{Serialize}(\mathbf{x}^i,\mathbf{y}^i, F) = \nonumber \\ 
&``\ f_1\ \text{is}\ \mathbf{x}^i_1.\ ... \ f_d\  \text{is}\  \mathbf{x}^i_d.\ \ \text{Answer:}\  \mathbf{y}^i\ ”
\end{align}

\vspace*{-0.12in}
\paragraph{Reasoning instruction.} To strengthen the reasoning abilities of LLMs, we supply targeted reasoning instructions (refer to the blue text in Fig.~\ref{fig:prompt_for_rule}). This guidance starts with an introductory phrase modeled after the chain-of-thought methodology~\cite{wei2022chain}. Subsequently, the reasoning is structured as a two-stage process. In the first stage, the LLM is prompted to hypothesize causal connections or tendencies between features and the task description independently of the provided examples, leveraging either general knowledge or common sense. This encourages the model to utilize its prior understanding relevant to the scenario. In the second stage, the LLM integrates the insights from the initial step along with example demonstrations to formulate classification rules. Segmenting reasoning in this way reduces the likelihood that the model relies on arbitrary associations from irrelevant attributes and encourages attention to the most pertinent features. To strike a balance between overly generic and excessively verbose rulesets—which could lead to underfitting or bloated prompts—we require extraction of a fixed quantity of rules (specifically 10)\footnote{Further analysis regarding the selection of rule quantity is available in Section~\ref{sec:analysis}.} for each class. Furthermore, during rule creation, the LLM references the feature types as detailed in the basic information section, ensuring compatibility and preventing rule type errors. 

\vspace*{-0.12in}
\paragraph{Response instruction.} For streamlined parsing and utilization of the generated rules, we furnish explicit directives regarding the response structure (illustrated by the yellow text in Fig.~\ref{fig:prompt_for_rule}). The prompt specifies the expected response layout for each answer class (i.e., response format), along with the prescribed structure each rule should conform to (i.e., rule format). These directions empower the LLM to choose formatting appropriate to each feature type. Where explicit formatting guidance is absent, the LLM may combine several rules with logical operators such as AND or OR. An illustrative example of the final output is available in Appendix~\ref{sec:outcome-rule}.

\subsection{Inferring Class Likelihood via Rules}
\label{Sec:training}

\paragraph{Rule-based feature extraction via parsing.} In this phase, we employ the rules identified in the earlier step to produce additional binary features. These features are generated separately for each class, specifying whether a given sample meets the criteria of that class's rule set (assigned values of '0' or '1'). Such features serve as indicators for estimating the likelihood that a sample pertains to a particular class. Because the LLM outputs rules in natural language, it becomes necessary to devise an automated parsing approach to facilitate feature extraction. Even though we attempt to regulate the structure and format during rule specification to simplify parsing, output from the LLM can still exhibit unpredictable syntactic modifications~\cite{kovavc2023large}. For example, the LLM could replace symbolic operators like ``$>$” or ``$<$” with textual expressions such as ``greater than" or ``less than," complicating automated text processing.

To mitigate difficulties in parsing irregular textual data, instead of relying on elaborate programmatic solutions, we turn to the LLM's capabilities directly. The LLM can reliably interpret semantic content even when syntax varies, and is adept at writing concise code when prompted (due to extensive exposure to code datasets during pretraining). To harness these strengths, we construct prompts containing the function name, descriptions of inputs and outputs, and the set of rules to be parsed, then feed these prompts into the LLM. Figures~\ref{fig:prompt_for_parsing} and ~\ref{sec:example_parsing}-\ref{sec:example_parse_outcome} in the Appendix provide illustrations of the prompt structure and its outcomes. The code produced by the LLM is executed using Python’s \texttt{exec()} function to carry out the data transformation.

\vspace*{-0.12in}
\paragraph{Estimating class probabilities.} After obtaining binary rule-based features for every class, a straightforward technique to quantify the likelihood of a sample's class membership is to aggregate the number of rules each class’s feature meets (i.e., by summing the corresponding binary values for each class). Yet, different rules may vary in their predictive utility, thus it is necessary to determine their relative importance through empirical learning on labeled data. To accomplish this, we employ traditional machine learning (ML) methods to learn the weight of each class's binary features. For instance, a simple bias-free linear model can be utilized. Let $\mathbf{z}_k^i \in \{0, 1\}^R$ represent the binary features obtained from sample $\mathbf{x}^i$ for class $k$, where $R$ stands for the number of rules associated with that class and $c$ denotes the number of classes. The likelihood for each class $\mathbf{p}^i$ is then computed as follows:

\begin{align}
\text{logit}_k^i &= \max(\mathbf{w}_k, 0) \cdot \mathbf{z}_k^i, \nonumber \\
\mathbf{p}^i &= \text{Softmax}({[\text{logit}_{1}^i, ..., \text{logit}_{c}^i]}). \label{eq:infer_prob}
\end{align}

In this context, $\mathbf{w}_k$ refers to the adjustable parameter associated with each feature in the linear model, quantifying its contribution to the prediction. By mapping these weights onto a non-negative subspace, we guarantee that the features produced by the LLM cannot negatively affect class probability during training. Subsequently, the resulting logit vector is transformed into class probabilities using the softmax operation. The parameters $\mathbf{w}_k$ are trained through a few-shot sample set by minimizing the cross-entropy objective.

\vspace*{-0.12in}
\paragraph{Ensembling with bagging.} In the final step, we iterate this full pipeline multiple times to assemble a collection of inference models, enabling prediction by ensemble. When provided with the set of learned weights across $T$ repetitions, $\{\mathbf{w}[t]\}_{t=1}^T$, prediction is performed by computing the mean probability for each class $\mathbf{p}^i$ using each weight vector $\mathbf{w}[t]$, as depicted in Eq.~\ref{eq:infer_prob}.

To foster diversity in the rules generated for the ensemble, a high sampling temperature is applied during LLM inference. Additionally, we permute the sequence of few-shot examples since empirical findings indicate that the order can sway LLM responses~\cite{lu2022fantastically}\footnote{A detailed examination of rule diversity is available in Appendix~\ref{sec:diversity}}. To further leverage prediction diversity for improved stability, bagging is used to choose a random subset of features or data instances in each run, a strategy that becomes especially beneficial when dealing with numerous training points or features. This ensemble methodology provides two notable benefits. Firstly, in cases where the LLM yields flawed rules for certain trials, the presence of accurate rules from other trials counteracts them, bolstering the system’s resilience. This relates to the notion of self-consistency in LLMs~\cite{wang2022self}, with frequently recurring rules across trials being more reliable. Secondly, the ensemble framework alleviates the issue of constrained prompt length in LLMs. When the tabular data surpasses the prompt’s input capacity, bagging allows for the reduction of data per trial, safeguarding model performance. Since a single rule set may not account for the interplay among all features, aggregating a variety of weak learners created over multiple trials mitigates the deficiencies associated with incomplete coverage in individual runs.

\section{Experiments}
We assess the effectiveness of \textsf{FeatLLM} across a range of tabular datasets and conduct multiple analyses to gain insight into our framework's operations. Owing to space limitations, supplementary experiments—including the exploration of alternative LLMs and qualitative assessments—are provided in the Appendix.

\subsection{Performance Evaluation}
\paragraph{Datasets.}
Thirteen datasets spanning binary and multiclass classification tasks are employed in our study: (1) Adult~\cite{asuncion2007uci}, which predicts if a person's annual income surpasses \$50,000; (2) Bank~\cite{moro2014data}, for forecasting the likelihood of a client subscribing to a term deposit; (3) Blood~\cite{yeh2009knowledge}, aimed at identifying whether blood donors return for future donations; (4) Car~\cite{kadra2021well}, which classifies vehicle quality; (5) Communities~\cite{misc_communities_and_crime_183}, used to estimate regional crime rates; (6) Credit-g~\cite{kadra2021well}, assessing individuals as good or bad credit risks; (7) Diabetes\footnote{\texttt{kaggle.com/datasets/uciml/pima-indians-diabetes-database}}, which predicts diabetes onset; (8) Heart\footnote{\texttt{kaggle.com/datasets/fedesoriano/heart-failure-prediction}}, for identifying coronary artery disease; and (9) Myocardial~\cite{misc_myocardial_infarction_complications_579}, predicting chronic heart failure status.

The experimental scope extends to contemporary and synthetic datasets that are either newly released or not covered during LLM pre-training: (10) Cultivars, which predicts soybean cultivar grain yield\footnote{\texttt{archive.ics.uci.edu/dataset/913}}; (11) NHANES, for classifying individuals’ age groups\footnote{\texttt{archive.ics.uci.edu/dataset/887}}; (12) Sequence-type, distinguishing sequences as arithmetic, geometric, Fibonacci, or Collatz types; and (13) Solution-mix, which determines if the concentration of a mixture comprising four solutions is above 0.5. The first two datasets became available to the UCI Machine Learning Repository post-September 2023, guaranteeing they are unseen in the LLM API (gpt-3.5-turbo-0613) utilized by our model. The latter two are artificially generated datasets, ensuring they are unfamiliar to the LLM, thus supporting a fair assessment.

Dataset sizes and their complexity, including attribute counts, are summarized in Table~\ref{Tab:basic-desc}. All datasets present explicit naming and documentation for each attribute. Datasets such as `Communities' and `Myocardial' are characterized by their high attribute dimensionality (over 100 features), which introduces obstacles given the restricted context length supported by current LLM architectures.

\vspace*{-0.12in}
\paragraph{Baselines.}
A total of nine baseline methods are used for comparative purposes. Three are established methods for tabular learning: (1) Logistic regression (LogReg), (2) XGBoost~\cite{chen2016xgboost}, and (3) RandomForest~\cite{ho1995random}. Two baselines operate under the assumption of accessible unlabeled data: (4) SCARF~\cite{bahri2022scarf}, and (5) STUNT~\cite{nam2023stunt}, although large unlabeled data is often challenging to obtain in practice. Another baseline, (6) TabPFN~\cite{hollmann2022tabpfn}, leverages pre-training on diverse synthetic distributions. The last three baselines are centered around LLMs: (7) In-context learning (In-context)~\cite{wei2022emergent}, (8) TABLET~\cite{slack2023tablet}, and (9) TabLLM~\cite{hegselmann2023tabllm}. The in-context learning approach simply places a few-shot training examples within the prompt, with no further tuning. TABLET augments the prompt with supplementary information derived from external classifiers, such as rule sets and prototypes, enhancing inference capabilities. TabLLM adopts parameter-efficient tuning methods (e.g., IA3~\cite{liu2022few}) for training with few-shot instances.

\vspace*{-0.12in}
\paragraph{Implementation details.}
\textsf{FeatLLM} leverages GPT-3.5 as its core LLM component, though its architecture is adaptable and does not depend on a specific LLM, as shown by results with alternate LLMs in Appendix~\ref{sec:backbones}. For inference, the LLM operates with a temperature of 0.5 and retains the API's default top-p setting of 1. We establish the ensemble size at 20 and set the rule extraction count to 10. Analysis of hyper-parameter choices can be found in Figure~\ref{fig:hyperparameter2} and Appendix~\ref{sec:hyper-parameter}. The linear model utilizes the Adam optimizer, configured with a 0.01 learning rate, and undergoes training for 200 epochs. To select the best epoch, we implement $k$-fold cross-validation. For baseline replication, GPT-3.5 is employed in in-context learning settings, while T0~\cite{sanh2022multitask} is used for TabLLM in accordance with its original configuration. It is important to note that TabLLM is limited to LLMs with publicly released checkpoints, thus GPT-3.5 cannot be applied in that scenario. Concerning classic machine learning algorithms including LogReg, XGBoost, and RandomForest, parameter tuning was accomplished using Grid-search in tandem with $k$-fold cross validation. The value of $k$ is selected as either 2 or 4, ensuring that every class has at least one representation in the training data. Comprehensive implementation details are available in Appendix~\ref{sec:baselines}.

\vspace*{-0.12in}
\paragraph{Results.}
Tables~\ref{tab:main1} and \ref{tab:main2} provide a comparative analysis of model performance across various datasets. Each experiment was conducted three times, employing different random partitions for training and test sets, and we report both the mean and standard deviation for the resulting AUC scores. In nearly all cases, \textsf{FeatLLM} either emerges as the leading model or occupies the second position relative to baseline competitors. Remarkably, our approach, even when constrained to just 4 shots, demonstrates performance on par with traditional tabular models trained with 32 shots (refer to Fig.~\ref{fig:compares}). While STUNT and TabPFN achieve substantial gains on particular datasets, their general improvement over established machine learning approaches remains limited. Furthermore, large language model (LLM) based frameworks—such as in-context learning, TABLET, or TabLLM—struggle to process tabular datasets with a high number of features. In contrast, our method, by incorporating bagging and ensembling, is able to handle high-dimensional datasets effectively, delivering robust results. It is also worth mentioning that the bagging and ensembling strategies we propose could potentially be applied to existing LLM-based models. Nevertheless, this would result in a twofold increase in inference computations per sample, significantly raising computational demands to an impractical level.

\subsection{Analysis \& Discussion}
\label{sec:analysis}
\paragraph{Ablation study.} To assess the influence of key components on overall performance, we conduct ablation studies, specifically examining: (1) \textsf{-Tuning}: eliminating the parameter tuning phase in the linear model and instead aggregating binary features to derive the logit; (2) \textsf{-Ensemble}: excluding the ensemble step; (3) \textsf{-Description}: removing the feature description element; and (4) \textsf{-Reasoning}: omitting the Step-1 reasoning instruction from rule generation.

Table~\ref{tab:ablation} shows the average AUC variation from each ablation, aggregated over all datasets. Any modification by ablation leads to a drop in performance. Interestingly, the advantage of parameter tuning increases with more shots, signifying that abundant data allows more precise rule weight estimation, amplifying the utility of tuning. Conversely, the positive contributions of feature descriptions and reasoning steps are more pronounced when fewer shots are available, highlighting the critical role of LLMs' prior knowledge under limited-data regimes. Ultimately, the proposed ensemble scheme persistently enhances results across all tested scenarios.

\vspace*{-0.12in}
\paragraph{Training \& inference time comparison.} We evaluate the efficiency of various methods in terms of both training and inference duration. Our experiments use the Adult dataset, where the training set consists of 16 examples. Unless specified otherwise, all procedures employ a single A100 GPU; however, TabLLM requires four A100 GPUs to facilitate model parallelism. Table~\ref{tab:cost} presents a detailed summary of these runtime evaluations. Significantly, \textsf{FeatLLM} achieves inference speeds on par with traditional baselines, maintaining a relatively low latency. Conversely, LLM-based techniques such as In-context learning, TABLET, and TabLLM incur substantially greater inference overhead. With respect to training, \textsf{FeatLLM}'s requirements are similar to those of other LLM-based methods, necessitating only 30 API requests—a modest demand that does not strain computational resources and allows for easy parallelization.

\vspace*{-0.12in}
\paragraph{Quality of features generated by \textsf{FeatLLM}.}
To assess the practical effectiveness of rule-based features produced by \textsf{FeatLLM}, we perform a straightforward test examining their informativeness. For this, we calculate the average AUROC between each new feature and the target variable, and report the proportion of features with AUROC exceeding 0.5 within each dataset. Results in Table~\ref{tab:quality} illustrate that most generated rules demonstrate clear relevance and offer valuable information for the prediction task (see the "Before tuning" column). Furthermore, when the linear model is adapted to the data, features with minimal predictive value (that is, those associated with weights below a specified threshold) are systematically excluded during tuning (see the "After tuning" column). The threshold for weight significance was established at 0.1, informed by the overall distribution of learned weights.

\vspace*{-0.12in}
\paragraph{Impact of introducing spurious correlations.} To evaluate how well different approaches can suppress the influence of irrelevant spurious correlations in situations where only a few labeled samples are available, we performed an experimental assessment. We employed the Heart dataset and augmented it by randomly appending columns from the Adult dataset, which are unrelated to the Heart task. Model performance was tracked as we gradually increased the number of these extraneous columns. All methods were provided an equal 8 labeled samples, and we did not utilize unlabeled data, thereby omitting algorithms such as SCARF and STUNT from this evaluation.

Figure~\ref{fig:spurious} illustrates how model performance fluctuates as more spurious correlations are introduced. Among the evaluated approaches, \textsf{FeatLLM} demonstrated the greatest resistance to performance decline induced by irrelevant correlations. This can be attributed to the framework’s reliance on prior knowledge to determine the relevance between the specific task and the features during rule extraction, which minimizes the likelihood of generating rules from unrelated columns and enhances robustness. Notably, rules built from irrelevant columns constituted only 1–2\% of all extracted rules. However, our analysis also indicates that \textsf{FeatLLM} may mistakenly infer spurious correlations or misidentify causal connections if columns derived from semantically related datasets are randomly introduced (refer to Appendix~\ref{sec:spurious-appendix}).

\vspace*{-0.12in}
\paragraph{Hyper-parameter analysis}
\textsf{FeatLLM} employs two principal hyper-parameters: the size of the ensemble and the number of rules per class derived in each LLM query. We systematically assessed their effects on model efficacy through a series of experiments. Our findings indicate that increasing ensemble size consistently enhances performance, but after a threshold (approximately 20 ensembles), the marginal improvements diminish and the results stabilize (refer to Fig.~\ref{fig:appendix-hyperparameter1} in Appendix). Regarding the rule count parameter, our evaluation revealed no significant variation in accuracy; however, choosing 10 rules strikes an optimal balance between prompt complexity and predictive accuracy (see Fig.~\ref{fig:hyperparameter2}). This outcome may be attributed to the expanded input dimensionality from additional rules, which, while informative, may induce overfitting in scenarios where training data are scarce.

\section{Conclusion}
We introduce \textsf{FeatLLM}, which establishes a new benchmark for few-shot tabular learning with LLMs. Unlike existing methods that employ LLMs to directly predict sample outputs, \textsf{FeatLLM} leverages LLMs to generate predictive criteria in the manner of feature engineering. The integration of rule creation and ensemble bagging means \textsf{FeatLLM} operates efficiently with minimal inference overhead, requires only API-level access—no separate model training—and avoids restrictions related to feature dimensionality. Demonstrating strong predictive accuracy, \textsf{FeatLLM} proves to be a highly effective choice for practical, data-efficient use of LLMs on real tabular datasets.

\textsf{FeatLLM} is optimized specifically for low-shot learning conditions. Looking ahead, we aim to expand its design to accommodate datasets with larger sample counts, diversify the kinds of features constructed beyond rules, and enhance interpretability, thereby broadening its applicability to various challenges.

\section{Broader Impact}
The framework, \textsf{FeatLLM}, is built to harness the pre-existing knowledge within LLMs for tabular learning tasks, with the intention of surpassing the capabilities of conventional supervised tabular methods. By prioritizing learning efficiency, this approach significantly addresses overfitting risk associated with scarce training data by incorporating prior knowledge. \textsf{FeatLLM} stands to benefit professionals in domains where label acquisition is resource-intensive or impractical (notably finance and healthcare), especially when LLMs are pretrained on sector-relevant texts. For broader adoption, it will be crucial to examine how societal biases and misinformation ingrained in LLM pretraining might influence outputs, particularly since such prior knowledge may assert itself over actual collected labeled data.

\end{document}